\chapter{Speculation and Deoptimization}\label{ch:speculation-deoptimization}

Optimized code bets on the future.
It assumes types stay stable, branches keep their direction, and call targets do not change.
When the bet holds, the code runs fast.
When the bet fails, the system must recover without the program noticing.
\term{Speculation} places those bets explicitly with checks called guards.
\term{Deoptimization} honors the fallback by rebuilding interpreter state when a guard fails.
This chapter makes both mechanisms precise enough to implement.
We place guards, record the state needed to leave, and trace one recovery by hand.

\section{Bets With Receipts}\label{sec:ch35-bets}

Every speculation has three parts: a fact assumed, a check that tests it, and a recovery that runs when the check fails.
The fact might be that a value is always an integer.
The check is a \term{guard}: a cheap test on the fast path that jumps to slow code on mismatch.
The recovery is a \term{deoptimization point} holding enough information to resume at a lower tier.
Figure~\ref{fig:ch35-guard-point} shows the shape.
Fast code runs guarded operations in sequence.
Each guard points to a shared exit stub.
The stub rebuilds interpreter state and resumes there.
Table~\ref{tab:ch35-speculations} lists common speculation kinds with their guard and fallback.
The table stays qualitative because guard costs depend on the target.
Type speculation guards tags and falls back to generic operations.
Branch speculation lays out the hot direction straight-line and falls back to the cold path.
Call speculation inlines the observed target and falls back to a full dispatch.
Alias and shape speculation assume a layout and fall back to generic access.
Each row follows the same discipline.
The assumed fact is written down, the check is emitted, and the fallback is tested before the optimization ships.

\begin{figure}[htbp]
\centering
\begin{tikzpicture}[node distance=1.0cm and 1.6cm]
\node[basic block] (fast1) {guarded op 1};
\node[basic block, below=of fast1] (fast2) {guarded op 2};
\node[basic block, below=of fast2] (fast3) {guarded op 3};
\node[basic block, right=of fast2] (stub) {exit stub\\ rebuild state};
\node[basic block, right=of stub] (interp) {interpreter\\ resume};
\draw[cfg edge] (fast1) -- (fast2);
\draw[cfg edge] (fast2) -- (fast3);
\draw[cfg edge] (fast1) -- (stub);
\draw[cfg edge] (fast2) -- (stub);
\draw[cfg edge] (fast3) -- (stub);
\draw[cfg edge] (stub) -- (interp);
\end{tikzpicture}
\caption{Each guarded operation on the fast path can reach a shared exit stub that rebuilds interpreter state, so no bet is placed without a tested way back.}
\label{fig:ch35-guard-point}
\end{figure}

\begin{table}[htbp]
\centering
\caption{Each speculation pairs an assumed fact with a guard and a tested fallback.}
\label{tab:ch35-speculations}
\begin{tabular}{lll}
\toprule
Assumption & Guard & Fallback \\
\midrule
Operand is integer & Tag test & Generic arithmetic \\
Branch goes one way & Direction test & Cold path \\
Call hits one target & Class test & Full dispatch \\
Object has one shape & Shape test & Generic access \\
\bottomrule
\end{tabular}
\end{table}

\section{What a Deoptimization Point Records}\label{sec:ch35-points}

A guard failure must reconstruct the exact state the interpreter would have held.
That requires a \term{frame descriptor} at each deoptimization point.
The descriptor maps compiled locations back to source-level values: locals, stack temporaries, and the bytecode index.
It also records which values need reboxing, re-tagging, or reloading from memory.
Algorithm~\ref{alg:ch35-deopt} gives the recovery procedure for one failed guard.
It assumes the descriptor for the failing point is complete and that the interpreter can resume at the recorded index.
It requires that no guarded effect before the failure becomes visible after recovery.
It copies each needed value through the descriptor, drops compiled-only temporaries, and resumes interpretation.
The ordering rule deserves emphasis.
Guards must fire before any effect the interpreter would not have produced yet.
Stores, calls with effects, and allocations constrain where guards may sit.
Compilers hoist pure checks freely but keep effectful operations behind their guards.
A frequent defect records a value still sitting in a reused register.
Liveness at the exit must reflect the interpreter, not the compiled code.
Another defect forgets a temporary that the interpreter needs two bytecodes later.
Descriptors therefore derive from interpreter state, computed forward, rather than from whatever happens to be live in machine code.

\begin{algorithm}[htbp]
\caption{Deoptimization after a failed guard}\label{alg:ch35-deopt}
\begin{pseudocode}
\Require failed guard $g$, complete descriptor $D(g)$, live compiled frame $F_c$; one logical interpreter frame
\Ensure interpreter frame $F_i$ resumed at the recorded index
\State $idx \gets$ bytecode index in $D(g)$
\State $F_i \gets$ fresh interpreter frame for the same method
\ForAll{$v$ live in the interpreter at $idx$}
\State $loc \gets$ compiled location of $v$ in $D(g)$
\State materialize $v$ in $F_i$ from $loc$ \Comment{rebox or reload as recorded}
\EndFor
\State discard $F_c$ temporaries not present in the interpreter
\State resume interpretation at $idx$ with $F_i$
\end{pseudocode}
\end{algorithm}

Materialization must preserve object identity as well as field contents: if two interpreter locals refer to one eliminated allocation, reconstruct one object and point both locals to it.
Keep source references visible to the collector while boxing or allocation can trigger collection, and retain the compiled frame until reconstruction succeeds.
Inlining extends the descriptor to a chain of logical frames; the single-frame algorithm is the base case.

\section{Worked Example: A Guard Fires Mid-Loop}\label{sec:ch35-worked}

We trace a loop that sums $1000$ inputs with one speculative integer add.
The optimizer guards each input's tag inside the loop; the next input can have a different tag, so that guard cannot simply be hoisted.
Iterations $0$ through $499$ see integers.
Iteration $500$ sees a string.
The compiled frame holds the running total unboxed in a register and the loop index in another.
The descriptor at the value guard records the bytecode index of the add, the boxed total, the index, and the current input slot.
Recovery under Algorithm~\ref{alg:ch35-deopt} proceeds step by step.
The guard at iteration $500$ fails on the string tag.
The stub reads the descriptor and boxes the running total accumulated so far.
Hand computation fixes the numbers: the integers $0$ through $499$ sum to $124750$.
The stub materializes total $124750$, index $500$, and the string input in a fresh interpreter frame.
It resumes at the recorded add index.
The interpreter performs generic addition for $124750$ and the string per language semantics.
If generic addition completes normally, the remaining $499$ iterations may continue interpreted or use later recompiled code; if it raises a language-defined error, that error is the required outcome.
No integer iteration is recomputed and no effect is duplicated.
Figure~\ref{fig:visual-deopt-map} makes the recovery mapping explicit.
The value $124750$ is the sum of inputs $0$ through $499$; input $500$ has been fetched but its addition has not executed.
Resuming after the add would omit an operation, while resuming before an already completed input read could duplicate an effect.
The descriptor must identify this semantic boundary as well as the storage locations.

\begin{figure}[htbp]
  \centering
  \begin{tikzpicture}[node distance=10mm and 22mm]
    \node[basic block,align=left] (compiled)
      {Compiled frame\\r4: 124750 (unboxed)\\r5: 500\\r6: current input};
    \node[basic block,align=left,right=of compiled] (interp)
      {Interpreter frame\\total: boxed 124750\\index: 500\\input: same value};
    \draw[cfg edge] (compiled) -- node[above,font=\small] {materialize} (interp);
    \node[basic block,below=of interp] (resume) {resume before add at pc 12};
    \draw[cfg edge] (interp) -- (resume);
    \node[font=\sffamily\small,align=center,below=of compiled]
      {The first 500 additions are complete.\\The failing addition has not run.};
  \end{tikzpicture}
  \caption{A deoptimization descriptor maps live compiled values into an interpreter state at the same semantic point, before the guarded addition.}
  \label{fig:visual-deopt-map}
\end{figure}


The listing below pins the descriptor contents for tests.
\begin{irlisting}
; deopt descriptor at speculative add
; bytecode index: 12 (add)
; total  <- reg r4, rebox as int
; index  <- reg r5, copy as int
; input  <- reg r6, rebox generic
; resume: interpreter at index 12
\end{irlisting}
A test replays this trace exactly and asserts the resumed total, the index, and the final result.

\section{Invalidation and Recompilation}\label{sec:ch35-invalidation}

Guards fail for two reasons: rare events and changed programs.
A rare event, like one string among thousands of integers, needs only local recovery.
Execution continues, counters note the event, and the compiled code stays installed.
Repeated failures mark the speculation as wrong and trigger recompilation with a wider assumption.
A changed program, like a class load that adds a method target, needs \term{invalidation}.
The system walks a dependency table from the changed entity to every compiled method that assumed the old fact.
Each affected method deoptimizes its live activations and discards or patches its code.
\begin{ccwarning}[Effects first]
Never move a guard below an effect it protects.
A store performed before a failing guard cannot be unperformed by the exit stub.
\end{ccwarning}
Test invalidation with a call site that first sees one target.
After inlining and guarding, a second target loads.
The test asserts that live activations of the old code resume correctly and that later calls use the new target.
The \enquote{load then call} sequence exercises the dependency between hierarchy changes and installed code; concurrent invalidation requires additional tests.

\section{Speculation Budgets and Failure Accounting}\label{sec:ch35-budgets}

Guards are cheap individually but unbounded in aggregate.
A method with hundreds of guards pays hundreds of tests on every pass, even when no bet ever fails.
A \term{speculation budget} caps how many active guards one method or loop may carry.
The optimizer ranks candidate speculations by expected gain and keeps only the affordable set.
Type stability at hot arithmetic usually outranks shape stability at cold property loads.
Call-site monomorphism usually outranks branch-direction layout for rarely taken forks.
\term{Failure accounting} turns guard exits into data for the next compilation.
Each deoptimization point counts its exits by cause: wrong type, new target, changed shape, or failed range.
One exit is noise.
Repeated exits with one cause mark the assumption as wrong and widen the next compilation.
Repeated exits with many causes mark the site as megamorphic and drop speculation there entirely.
Consider a call site that first sees target $A$ ninety-eight times, then $B$ once, then $C$ once.
A single-target speculation with a fallback covers this shape well.
If later runs show $A$, $B$, and $C$ in equal thirds, the accounting promotes the site and the next compile uses a small dispatch instead of a guard.
Test budgets by compiling a method with more candidates than its allowance and asserting the kept set matches the ranking.
Test accounting by replaying a recorded exit mix and asserting the recompilation decision.
The \enquote{decide from counts} rule keeps speculation policy reproducible and reviewable.

\section{Implementation Notes and Complexity}\label{sec:ch35-implementation}

Each guard costs one test and one rarely-taken branch, typically $O(1)$.
Descriptor size grows with the live interpreter values at each point.
Sharing descriptors across nearby points and compressing locations keeps metadata small.
Deoptimization itself costs $O(V)$ where $V$ is the live-value count at the exit, plus frame allocation.
Invalidation costs $O(A)$ where $A$ is the affected-method count found through dependencies.
Production systems tier their response: count single failures, recompile after a few, and invalidate broadly only on program change.
They also limit speculation depth so nested inlines do not multiply descriptors.
A common pitfall guards too much and spends the speedup on tests.
Another pitfall records descriptors lazily and discovers gaps only when a guard fires in production.
We record eagerly, test exits with injected failures, and measure guard overhead separately from body speed.
Speculation here connects to tiering in the previous chapter and to vector guards in the next.
The descriptor discipline also previews the testing chapter: every exit is a contract with a replayable test.
% TODO(cite): need a primary reference for deoptimization-point liveness rules.

\section{Exercises}\label{sec:ch35-exercises}

\begin{enumerate}
\item A loop starts with total zero and adds integer input $i$ at iteration $i$ for $0\le i<600$; input $600$ is a float. Place the guards in this $800$-iteration loop, compute the total before the failed guard, and list the descriptor entries needed to resume.\label{ex:ch35-guard-math}
\item Distinguish a rare-event guard failure from a program-change invalidation, stating for each whether code stays installed and what counters should record.\label{ex:ch35-rare-vs-change}
\item Write a test that injects a guard failure at a chosen point, listing the expected interpreter values after recovery and the assertion that catches a dropped temporary.\label{ex:ch35-inject-test}
\end{enumerate}

\section{Further Reading}\label{sec:ch35-reading}

Appel presents the analysis and transformation ideas that speculation builds upon~\cite{appel-1998-modern}.
Morgan discusses profile-driven code shape and its risks~\cite{morgan-1998-optimizing}.
Safonov treats assumption changes as trustworthiness obligations~\cite{safonov-2010-trustworthy}.
The C standard draft defines the observable effects that constrain guard placement~\cite{wg14-2011-c11}.
